\chapter{Exercício 3}
\label{cha:exercise_3}

\section{Cadeia de certificados do site example.org}

Ao aceder ao site \url{https://www.example.org}, verificamos que a cadeia de certificados é composta por três elementos:

\begin{enumerate}
  \item \textbf{Certificado do servidor:} \texttt{*.example.org}, emitido pela \textit{Internet Corporation for Assigned Names and Numbers (ICANN)}.
  \item \textbf{Certificado intermédio:} \texttt{DigiCert Global G3 TLS ECC SHA384 2020 CA1}, emitido pela autoridade de certificação DigiCert.
  \item \textbf{Certificado raiz:} \texttt{DigiCert Global Root G3}, autoassinado e presente no armazenamento de autoridades de certificação confiáveis dos sistemas operativos e navegadores.
\end{enumerate}

A cadeia garante que o certificado do servidor é confiável, desde que a autoridade raiz seja confiável no sistema do cliente.

\section{Algoritmo de assinatura do certificado}

Usando o comando OpenSSL:

\begin{verbatim}
openssl s_client -connect example.org:443 -showcerts
\end{verbatim}

é possível visualizar o campo \texttt{sigalg}, que indica o algoritmo usado para assinar o certificado:

\begin{verbatim}
sigalg: ecdsa-with-SHA384
\end{verbatim}

Portanto, o certificado foi assinado utilizando o algoritmo \textbf{ECDSA (Elliptic Curve Digital Signature Algorithm)} em conjunto com a função de hash \textbf{SHA-384}.
Isto garante a integridade do certificado e permite verificar que não foi alterado após a assinatura pela autoridade certificadora.